\chapter{Development Trace: Historical Pre-Leakage Runs}
\label{app:dev-trace}

This appendix preserves the development trace of the re-ranker work: the
sequence of early experiments (Experiments~1 to~7) that shaped the final
methodology but predate the leakage fix and are therefore \emph{not}
live-reproducible from the PROTEA app. They were run on the retired 22-feature
\texttt{bench-v1-K5} schema over the GOA~220~$\to$~229 evaluation split, before
the anc2vec look-up leakage was identified and patched and before the
\texttt{cafaeval} (\texttt{prop=fill, norm=cafa}) protocol was adopted for
reporting. Their absolute $F_{\max}$ values are not directly comparable to the
leakage-clean numbers in the body of \cref{chap:evaluation}, and they carry no
live \texttt{/benchmark} or \texttt{/evaluation} source. They are retained here
as an auditable record of the design path, including its dead-ends, and provide
the empirical trace behind \textbf{RQ4}.

Every result in this appendix is provenance-anchored to the frozen lab CSV and
commit cited in the corresponding figure caption rather than to the live app.
The leakage-clean, live-reproducible studies that constitute the primary
results (the 56-feature study and the binary-label champion with its multi-PLM
grid) are reported in \cref{sec:experiments-results}.

% ── Experiment 1 ──────────────────────────────────────────────
\section[Experiment~1: Effect of k (Neighbours Per Query)]{Experiment~1: Effect of $k$ (Neighbours Per Query)}
\label{subsec:exp1-k}

The first experiment establishes the optimal number of nearest neighbours $k$
for annotation transfer. Four values are tested: $k \in \{5, 10, 20, 50\}$,
all using \texttt{aspect\_separated\_knn=true} and the baseline scoring
$c = 1 - d/2$.

\begin{table}[htbp]
  \centering
  \caption[Fmax vs.\ number of neighbours k]{$F_{\max}$ vs.\ number of neighbours $k$.}
  \label{tab:results-k-sweep}
  \resizebox{\textwidth}{!}{%
  \begin{tabular}{l ccc ccc ccc}
    \toprule
    & \multicolumn{3}{c}{\textbf{NK}} & \multicolumn{3}{c}{\textbf{LK}} & \multicolumn{3}{c}{\textbf{PK}} \\
    \cmidrule(lr){2-4} \cmidrule(lr){5-7} \cmidrule(lr){8-10}
    $k$ & BPO & MFO & CCO & BPO & MFO & CCO & BPO & MFO & CCO \\
    \midrule
    \textbf{5}  & \textbf{0.412} & \textbf{0.590} & \textbf{0.668} & \textbf{0.467} & \textbf{0.558} & \textbf{0.676} & \textbf{0.187} & \textbf{0.278} & \textbf{0.325} \\
    10 & 0.400 & 0.574 & 0.656 & 0.458 & 0.537 & 0.663 & 0.177 & 0.272 & 0.317 \\
    20 & 0.396 & 0.564 & 0.649 & 0.454 & 0.528 & 0.654 & 0.173 & 0.269 & 0.313 \\
    50 & 0.396 & 0.555 & 0.646 & 0.452 & 0.523 & 0.651 & 0.173 & 0.269 & 0.312 \\
    \bottomrule
  \end{tabular}
  }
\end{table}

$k{=}5$ is optimal across all categories and aspects. Increasing $k$ degrades
performance monotonically: additional neighbours introduce noise without
meaningful recall gains. This is consistent with the concentration of function
in the nearest embedding-space neighbours. All subsequent experiments use
$k{=}5$.

\begin{figure}[htbp]
  \centering
  % ============================================================
%  Figure: Fmax as a function of k (Experiment 1)
%  Data source: figures/data/fmax_vs_k.csv
%  Generated by: figures/fig_fmax_vs_k.py
% ============================================================
\begin{tikzpicture}
  \begin{groupplot}[
    group style={
      group size=3 by 1,
      horizontal sep=10mm,
      ylabels at=edge left,
    },
    protea plot,
    xmode=log,
    log basis x=2,
    xtick={5,10,20,50},
    xticklabels={5,10,20,50},
    xlabel={$k$ (neighbours)},
    ylabel={$F_{\max}$},
    ymin=0.15, ymax=0.72,
    grid=both,
    legend columns=3,
  ]
    % -- NK panel
    %  legend to name is attached to a single panel only. Keeping it in the
    %  group-wide options makes every panel re-emit the same label, which
    %  triggers "multiply defined label" warnings.
    \nextgroupplot[title={NK (No-Knowledge)}, legend to name=fmaxk-legend]
      \addplot coordinates {(5,0.412) (10,0.400) (20,0.396) (50,0.396)};
      \addlegendentry{BPO}
      \addplot coordinates {(5,0.590) (10,0.574) (20,0.564) (50,0.555)};
      \addlegendentry{MFO}
      \addplot coordinates {(5,0.668) (10,0.656) (20,0.649) (50,0.646)};
      \addlegendentry{CCO}

    % -- LK panel
    \nextgroupplot[title={LK (Limited-Knowledge)}]
      \addplot coordinates {(5,0.467) (10,0.458) (20,0.454) (50,0.452)};
      \addlegendentry{BPO}
      \addplot coordinates {(5,0.558) (10,0.537) (20,0.528) (50,0.523)};
      \addlegendentry{MFO}
      \addplot coordinates {(5,0.676) (10,0.663) (20,0.654) (50,0.651)};
      \addlegendentry{CCO}

    % -- PK panel
    \nextgroupplot[title={PK (Partial-Knowledge)}]
      \addplot coordinates {(5,0.187) (10,0.177) (20,0.173) (50,0.173)};
      \addlegendentry{BPO}
      \addplot coordinates {(5,0.278) (10,0.272) (20,0.269) (50,0.269)};
      \addlegendentry{MFO}
      \addplot coordinates {(5,0.325) (10,0.317) (20,0.313) (50,0.312)};
      \addlegendentry{CCO}
  \end{groupplot}

  % Centered legend below the group
  \node[anchor=north] at ($(group c2r1.south) + (0,-7mm)$)
    {\pgfplotslegendfromname{fmaxk-legend}};
\end{tikzpicture}

  \caption[Fmax as a function of k]{%
    $F_{\max}$ as a function of the number of neighbours $k$ for each evaluation
    category and GO aspect. Performance degrades monotonically as $k$ grows,
    confirming that functional signal is concentrated in the very nearest
    neighbours and that aggregating over more candidates merely dilutes the vote.
    Note the log-scale on the $x$ axis.
    Source: \texttt{figures/fig\_fmax\_vs\_k.py} @ \texttt{21d74f7}
    (data: \texttt{figures/data/fmax\_vs\_k.csv}).%
  }
  \label{fig:fmax-vs-k}
\end{figure}

% ── Experiment 2 ──────────────────────────────────────────────
\section{Experiment~2: Aspect-Separated vs.\ Unified KNN}
\label{subsec:exp2-aspect}

In aspect-separated mode, three independent \gls{knn} indices are built, one
per GO namespace, using only reference proteins annotated in that namespace.
This guarantees coverage of all three aspects but may dilute the signal for
well-represented namespaces.

\begin{table}[htbp]
  \centering
  \caption[Fmax: aspect-separated vs.\ unified KNN]{$F_{\max}$: aspect-separated vs.\ unified KNN ($k{=}5$).}
  \label{tab:results-aspect-sep}
  \resizebox{\textwidth}{!}{%
  \begin{tabular}{l ccc ccc ccc}
    \toprule
    & \multicolumn{3}{c}{\textbf{NK}} & \multicolumn{3}{c}{\textbf{LK}} & \multicolumn{3}{c}{\textbf{PK}} \\
    \cmidrule(lr){2-4} \cmidrule(lr){5-7} \cmidrule(lr){8-10}
    & BPO & MFO & CCO & BPO & MFO & CCO & BPO & MFO & CCO \\
    \midrule
    Separated  & 0.412 & 0.590 & 0.668 & 0.467 & 0.558 & 0.676 & 0.187 & 0.278 & 0.325 \\
    Unified    & 0.410 & 0.595 & 0.666 & 0.471 & 0.569 & 0.675 & 0.188 & 0.279 & 0.325 \\
    \bottomrule
  \end{tabular}
  }
\end{table}

Differences are minimal ($\leq 0.011$). The unified index slightly favours \gls{mfo}
while the separated index slightly favours \gls{bpo}. We retain
\texttt{aspect\_separated=true} for all subsequent experiments as it provides
uniform aspect coverage without measurable degradation.

% ── Experiment 3 ──────────────────────────────────────────────
\section{Experiment~3: Heuristic Scoring Configurations}
\label{subsec:exp3-scoring}

This experiment tests whether combining embedding distance with alignment
statistics and evidence codes improves scoring without machine learning.
A single prediction set with \texttt{compute\_alignments=true} and
\texttt{compute\_taxonomy=true} is generated; five scoring configurations
are applied post-hoc via different weight vectors.

\begin{table}[htbp]
  \centering
  \caption[Fmax under five heuristic scoring configurations]{$F_{\max}$ under five heuristic scoring configurations.}
  \label{tab:results-scoring}
  \resizebox{\textwidth}{!}{%
  \begin{tabular}{l ccc ccc ccc}
    \toprule
    & \multicolumn{3}{c}{\textbf{NK}} & \multicolumn{3}{c}{\textbf{LK}} & \multicolumn{3}{c}{\textbf{PK}} \\
    \cmidrule(lr){2-4} \cmidrule(lr){5-7} \cmidrule(lr){8-10}
    \textbf{Config} & BPO & MFO & CCO & BPO & MFO & CCO & BPO & MFO & CCO \\
    \midrule
    embedding\_only     & 0.412 & 0.590 & 0.668 & 0.467 & 0.558 & 0.675 & 0.187 & 0.278 & 0.325 \\
    \textbf{alignment\_weighted} & \textbf{0.428} & \textbf{0.611} & \textbf{0.683} & \textbf{0.500} & \textbf{0.598} & \textbf{0.699} & \textbf{0.201} & \textbf{0.285} & \textbf{0.337} \\
    evidence\_primary   & 0.362 & 0.558 & 0.638 & 0.412 & 0.540 & 0.642 & 0.165 & 0.268 & 0.308 \\
    emb+evidence        & 0.352 & 0.531 & 0.618 & 0.387 & 0.517 & 0.626 & 0.162 & 0.250 & 0.300 \\
    composite           & 0.364 & 0.560 & 0.639 & 0.412 & 0.542 & 0.642 & 0.167 & 0.267 & 0.307 \\
    \bottomrule
  \end{tabular}
  }
\end{table}

The \texttt{alignment\_weighted} configuration (embedding~$=0.5$, NW
identity~$=0.3$, SW identity~$=0.2$) improves the baseline by 1.5 to 4\%
$F_{\max}$ across all cells. Configurations that incorporate evidence code
weights \emph{degrade} performance: the evidence code signal conflicts with
IA weighting used in evaluation. This is an early indication that pairwise
alignment statistics provide a complementary signal to embedding distance,
a question revisited under the leakage-clean protocol in
\cref{sec:experiments-results}.

% ── Experiment 4 ──────────────────────────────────────────────
\section{Experiment~4: LightGBM Re-Ranker, Per-Aspect Models}
\label{subsec:exp4-reranker-pa}

A LightGBM binary classifier is trained on 12 temporal holdout splits
(GOA~160$\to$165 through GOA~215$\to$220) and evaluated on the test split
(GOA~220$\to$229). Nine models are trained (NK/LK/PK $\times$ BPO/MFO/CCO).
Two variants are tested: no class balancing, and subsampling negatives to a
10:1 ratio (\texttt{neg\_pos\_ratio=10}).

\begin{table}[htbp]
  \centering
  \caption[Fmax: LightGBM per-aspect re-ranker]{$F_{\max}$: LightGBM per-aspect re-ranker (Exp.~4).}
  \label{tab:results-reranker-pa}
  \resizebox{\textwidth}{!}{%
  \begin{tabular}{l ccc ccc ccc}
    \toprule
    & \multicolumn{3}{c}{\textbf{NK}} & \multicolumn{3}{c}{\textbf{LK}} & \multicolumn{3}{c}{\textbf{PK}} \\
    \cmidrule(lr){2-4} \cmidrule(lr){5-7} \cmidrule(lr){8-10}
    \textbf{Method} & BPO & MFO & CCO & BPO & MFO & CCO & BPO & MFO & CCO \\
    \midrule
    Baseline           & 0.412 & 0.590 & 0.668 & 0.467 & 0.558 & 0.675 & 0.187 & 0.278 & 0.325 \\
    Alignment weighted & 0.428 & 0.611 & 0.683 & 0.500 & 0.598 & 0.699 & 0.201 & 0.285 & 0.337 \\
    PA (no balance)    & 0.384 & 0.584 & \textbf{0.695} & 0.447 & 0.482 & \textbf{0.713} & 0.201 & 0.284 & 0.335 \\
    PA (balanced)      & 0.408 & 0.577 & 0.687 & 0.478 & 0.506 & 0.711 & 0.201 & 0.298 & 0.332 \\
    \bottomrule
  \end{tabular}
  }
\end{table}

The per-aspect re-ranker improves CCO substantially (+2 to 4\% over baseline) but
\emph{degrades} MFO by 3 to 9\% relative to the heuristic. Six of nine models
suffer from early stopping at iteration~1 due to extreme class imbalance
(positive rate $<0.2\%$). Class balancing corrects the BPO collapse but
introduces a different trade-off. Neither variant surpasses the heuristic
overall.

% ── Experiment 5 ──────────────────────────────────────────────
\section{Experiment~5: LightGBM Re-Ranker, Per-Category Models with IA Weighting}
\label{subsec:exp5-reranker-pc}

Experiment~5 addresses three weaknesses of the per-aspect re-ranker (Exp.~4):
\begin{enumerate}
  \item \textbf{Per-category models} (NK, LK, PK) instead of per-aspect
    (NK-BPO, NK-MFO, \ldots), giving each model ${\sim}3{\times}$ more training data.
  \item \textbf{IA sample weighting}: each training example is weighted by the
    Information Accretion of its GO term, aligning the loss with the evaluation metric.
  \item Hyperparameter refinement: learning rate $0.05 \to 0.01$,
    max rounds $300 \to 1000$, early stopping patience $30 \to 50$.
\end{enumerate}

\begin{table}[htbp]
  \centering
  \caption[Fmax: LightGBM per-category re-ranker]{$F_{\max}$: LightGBM per-category re-ranker with IA weighting (Exp.~5).}
  \label{tab:results-reranker-pc}
  \resizebox{\textwidth}{!}{%
  \begin{tabular}{l ccc ccc ccc}
    \toprule
    & \multicolumn{3}{c}{\textbf{NK}} & \multicolumn{3}{c}{\textbf{LK}} & \multicolumn{3}{c}{\textbf{PK}} \\
    \cmidrule(lr){2-4} \cmidrule(lr){5-7} \cmidrule(lr){8-10}
    \textbf{Method} & BPO & MFO & CCO & BPO & MFO & CCO & BPO & MFO & CCO \\
    \midrule
    Baseline           & 0.412 & 0.590 & 0.668 & 0.467 & 0.558 & 0.675 & 0.187 & 0.278 & 0.325 \\
    Alignment weighted & 0.428 & 0.611 & 0.683 & 0.500 & 0.598 & 0.699 & 0.201 & 0.285 & 0.337 \\
    PC (full)          & 0.425 & 0.607 & 0.689 & 0.486 & 0.575 & 0.707 & 0.199 & 0.297 & 0.335 \\
    \bottomrule
  \end{tabular}
  }
\end{table}

The per-category re-ranker is substantially more robust than the per-aspect
configuration: MFO no longer collapses (0.607 vs.\ 0.577) and performance is
consistent across all cells. However, \texttt{alignment\_weighted} still leads
in BPO and MFO. Post-hoc analysis of feature importance revealed that alignment
and taxonomy features had \textbf{zero importance} in all three models: these
features were set to \texttt{NULL} during training data generation, meaning the
model could only learn from embedding distance, aggregation features, and
categorical signals.

% ── Experiment 6 ──────────────────────────────────────────────
\section{Experiment~6: LightGBM Re-Ranker, Full Feature Set}
\label{subsec:exp6-reranker-ff}

Experiment~6 addresses the missing-feature problem identified in Exp.~5 by
computing pairwise alignment (Needleman-Wunsch and Smith-Waterman via
parasail/BLOSUM62), taxonomic distance (via the NCBI taxonomy tree), and
anc2vec GO-embedding cosine similarities
for every (query, reference) pair during training data generation. This gives the
model access to all 56 features, the same set available at inference time.

\textbf{Note on feature schema at time of writing.} Subsequent development
(T-RES.1b, 2026-05-12) added four lineage features
(\texttt{lineage\_is\_ancestor\_of\_known}, \texttt{lineage\_is\_descendant\_of\_known},
\texttt{lineage\_ancestor\_of\_count}, \texttt{lineage\_descendant\_of\_count}) to
the canonical \texttt{FeatureRegistry}, bringing the schema from 52 to 56 columns.
The 56-feature study described in Experiment~8 (\cref{subsec:exp8-reranker-56f}) was run
against the current 56-column schema.

Training uses the same 12 temporal holdout splits as Exp.~5, with identical
hyperparameters. The feature dataset is the leakage-free
\texttt{bench-v1-K5-filtered} snapshot (evaluation pair GOA~220$\to$230),
which removes anc2vec look-up leakage present in the original
\texttt{bench-v1-K5} dataset. Nine per-cell models are trained
(NK/LK/PK $\times$ BPO/MFO/CCO) and evaluated with the lab
protein-averaged $F_{\max}$ metric (no prediction-side propagation); three
random seeds (42, 7, 137) are run and averaged.

\begin{quote}
\textbf{Note.} The preceding experiments (Exp.~1 to 5) were run on the
pre-leakage-fix \texttt{bench-v1-K5} dataset with the older 22-feature schema.
Their absolute $F_{\max}$ values are not directly comparable to the full-feature
re-ranker numbers reported here. Under the current protocol, a controlled
cross-method comparison requires re-running the baseline and heuristic
configurations on a shared leakage-clean SELECT-window split; that re-run is
reported as future work, and the earlier rows are retained for directional
reference only.
\end{quote}

\begin{table}[htbp]
  \centering
  \caption[Fmax: complete comparison of all methods]{$F_{\max}$: complete comparison of all methods.}
  \label{tab:results-all}
  \resizebox{\textwidth}{!}{%
  \begin{tabular}{l ccc ccc ccc}
    \toprule
    & \multicolumn{3}{c}{\textbf{NK}} & \multicolumn{3}{c}{\textbf{LK}} & \multicolumn{3}{c}{\textbf{PK}} \\
    \cmidrule(lr){2-4} \cmidrule(lr){5-7} \cmidrule(lr){8-10}
    \textbf{Method} & BPO & MFO & CCO & BPO & MFO & CCO & BPO & MFO & CCO \\
    \midrule
    Baseline (emb.\ only)$^\dagger$ & 0.412 & 0.590 & 0.668 & 0.467 & 0.558 & 0.675 & 0.187 & 0.278 & 0.325 \\
    Alignment weighted$^\dagger$    & 0.428 & 0.611 & 0.683 & 0.500 & 0.598 & 0.699 & 0.201 & 0.285 & 0.337 \\
    Per-aspect (bal.)$^\dagger$     & 0.408 & 0.577 & 0.687 & 0.478 & 0.506 & 0.711 & 0.201 & 0.298 & 0.332 \\
    Per-category$^\dagger$          & 0.425 & 0.607 & 0.689 & 0.486 & 0.575 & 0.707 & 0.199 & 0.297 & 0.335 \\
    \textbf{Full-feature} & \textbf{0.599} & \textbf{0.646} & \textbf{0.775} & \textbf{0.620} & \textbf{0.606} & \textbf{0.794} & \textbf{0.165} & \textbf{0.293} & \textbf{0.346} \\
    \bottomrule
  \end{tabular}
  }
\end{table}

\noindent$^\dagger$~Numbers from pre-leakage-fix \texttt{bench-v1-K5} dataset
(22-feature schema, GOA~220$\to$229 eval split); not directly comparable with
the full-feature row, which uses \texttt{bench-v1-K5-filtered}
(56-feature schema, GOA~220$\to$230 eval split, mean of 3 seeds).

The full-feature re-ranker trained on the leakage-free feature set achieves strong
$F_{\max}$ across all nine category-aspect cells. The GOA~220$\to$230
evaluation window (versus 220$\to$229 for earlier rows) adds one additional
GOA release to the test delta; this difference is expected to be small relative
to the dataset size. Direct cell-by-cell comparison with the pre-leakage-fix
rows is not valid; a controlled leakage-free re-run of the earlier methods on
the shared SELECT-window split is future work.

\begin{figure}[htbp]
  \centering
  % ============================================================
%  Figure: Re-ranker stage progression (Experiments 4 to 6)
%  Data source: figures/data/reranker_progression.csv
%  Generated by: figures/fig_reranker_progression.py
%  Stages: Baseline, AlnWeighted, PerAspect, PerCategory, FullFeature
%  NOTE: Baseline/AlnWeighted/PerAspect/PerCategory use the
%        pre-leakage-fix bench-v1-K5 dataset (22-feature schema,
%        GOA 220->229).
%        FullFeature uses leakage-free bench-v1-K5-filtered
%        (52-feature schema, GOA 220->230, mean of 3 seeds).
%        The earlier stages are directional reference only.
% ============================================================
\resizebox{\textwidth}{!}{%
\begin{tikzpicture}
  \begin{groupplot}[
    group style={
      group size=3 by 1,
      horizontal sep=12mm,
      ylabels at=edge left,
    },
    protea plot,
    width=6.0cm, height=5.0cm,
    enlarge x limits=0.18,
    symbolic x coords={Baseline, AlnWeighted, PerAspect, PerCategory, FullFeature},
    xtick=data,
    xticklabel style={font=\tiny, rotate=35, anchor=east, yshift=0.5mm},
    ymin=0.10, ymax=0.85,
    ylabel={$F_{\max}$},
    grid=both,
    legend columns=3,
    every axis plot/.append style={fill opacity=0.85},
  ]

    % -- NK panel
    \nextgroupplot[title={NK}, ybar, bar width=2.7pt, legend to name=reranker-legend]
      \addplot+[fill=proteagreen, draw=proteagreen!60!black]
        coordinates {(Baseline,0.412) (AlnWeighted,0.428) (PerAspect,0.408) (PerCategory,0.425) (FullFeature,0.599)};
      \addlegendentry{BPO}
      \addplot+[fill=linkblue, draw=linkblue!60!black]
        coordinates {(Baseline,0.590) (AlnWeighted,0.611) (PerAspect,0.577) (PerCategory,0.607) (FullFeature,0.646)};
      \addlegendentry{MFO}
      \addplot+[fill=BurntOrange, draw=BurntOrange!60!black]
        coordinates {(Baseline,0.668) (AlnWeighted,0.683) (PerAspect,0.687) (PerCategory,0.689) (FullFeature,0.775)};
      \addlegendentry{CCO}

    % -- LK panel
    \nextgroupplot[title={LK}, ybar, bar width=2.7pt]
      \addplot+[fill=proteagreen, draw=proteagreen!60!black]
        coordinates {(Baseline,0.467) (AlnWeighted,0.500) (PerAspect,0.478) (PerCategory,0.486) (FullFeature,0.620)};
      \addplot+[fill=linkblue, draw=linkblue!60!black]
        coordinates {(Baseline,0.558) (AlnWeighted,0.598) (PerAspect,0.506) (PerCategory,0.575) (FullFeature,0.606)};
      \addplot+[fill=BurntOrange, draw=BurntOrange!60!black]
        coordinates {(Baseline,0.675) (AlnWeighted,0.699) (PerAspect,0.711) (PerCategory,0.707) (FullFeature,0.794)};

    % -- PK panel
    \nextgroupplot[title={PK}, ymin=0.10, ymax=0.85, ybar, bar width=2.7pt]
      \addplot+[fill=proteagreen, draw=proteagreen!60!black]
        coordinates {(Baseline,0.187) (AlnWeighted,0.201) (PerAspect,0.201) (PerCategory,0.199) (FullFeature,0.165)};
      \addplot+[fill=linkblue, draw=linkblue!60!black]
        coordinates {(Baseline,0.278) (AlnWeighted,0.285) (PerAspect,0.298) (PerCategory,0.297) (FullFeature,0.293)};
      \addplot+[fill=BurntOrange, draw=BurntOrange!60!black]
        coordinates {(Baseline,0.325) (AlnWeighted,0.337) (PerAspect,0.332) (PerCategory,0.335) (FullFeature,0.346)};
  \end{groupplot}

  \node[anchor=north] at ($(group c2r1.south) + (0,-12mm)$)
    {\pgfplotslegendfromname{reranker-legend}};
\end{tikzpicture}%
}

  \caption[Re-ranker progression across methods]{%
    Method progression faceted by category (NK / LK / PK) with one bar per GO
    aspect. The full-feature bar reflects the leakage-free \texttt{bench-v1-K5-filtered}
    dataset (56 features, mean of 3 seeds); all other bars use the
    pre-leakage-fix \texttt{bench-v1-K5} dataset and are included for
    directional reference only. A controlled leakage-free re-run of the earlier
    stages on the shared SELECT-window split is future work.
    Source: \texttt{figures/fig\_reranker\_progression.py} @ \texttt{21d74f7}
    (data: \texttt{figures/data/reranker\_progression.csv}).%
  }
  \label{fig:reranker-progression}
\end{figure}

\subsection{Feature Importance Analysis}

\Cref{tab:feature-importance} shows the top-10 features by LightGBM gain
for each category, aggregated over the three per-aspect models (BPO, MFO, CCO),
seed 42, leakage-free dataset.

\begin{table}[htbp]
  \centering
  \caption{Top-10 features by LightGBM gain for full-feature per-cell models (seed~42,
    aggregated across BPO + MFO + CCO within each category).}
  \label{tab:feature-importance}
  \resizebox{\textwidth}{!}{%
  \begin{tabular}{clr clr clr}
    \toprule
    \multicolumn{3}{c}{\textbf{NK}} & \multicolumn{3}{c}{\textbf{LK}} & \multicolumn{3}{c}{\textbf{PK}} \\
    \cmidrule(lr){1-3} \cmidrule(lr){4-6} \cmidrule(lr){7-9}
    \# & Feature & Gain & \# & Feature & Gain & \# & Feature & Gain \\
    \midrule
    1 & nb\_vote\_frac   & 155K & 1 & nb\_vote\_frac   & 98K  & 1 & go\_term\_freq       & 410K \\
    2 & k\_position      & 58K  & 2 & evidence\_code   & 45K  & 2 & anc2vec\_q\_maxcos   & 233K \\
    3 & evidence\_code   & 48K  & 3 & k\_position      & 38K  & 3 & nb\_vote\_frac       & 166K \\
    4 & go\_term\_freq   & 45K  & 4 & go\_term\_freq   & 33K  & 4 & nb\_min\_dist        & 133K \\
    5 & vote\_count      & 28K  & 5 & anc2vec\_n\_cos  & 18K  & 5 & qualifier            & 85K  \\
    6 & anc2vec\_n\_cos  & 27K  & 6 & vote\_count      & 17K  & 6 & anc2vec\_n\_cos      & 83K  \\
    7 & ref\_ann\_dens.  & 17K  & 7 & anc2vec\_q\_maxcos& 13K & 7 & evidence\_code       & 64K  \\
    8 & emb\_pca\_q\_5   & 5K   & 8 & ref\_ann\_dens.  & 9K   & 8 & anc2vec\_q\_cos      & 55K  \\
    9 & emb\_pca\_q\_0   & 5K   & 9 & anc2vec\_q\_count& 7K   & 9 & vote\_count          & 53K  \\
   10 & anc2vec\_n\_maxcos& 5K  & 10 & anc2vec\_q\_cos & 7K   & 10 & anc2vec\_q\_count   & 50K  \\
    \bottomrule
  \end{tabular}
  }
\end{table}

\begin{figure}[htbp]
  \centering
  % ============================================================
%  Figure: LightGBM feature importance for v3 re-ranker
%  Data source: Table 6.5 (chapters/06_evaluation.tex)
%  Gain values are in thousands.
%  Symbolic ids are ASCII; pretty labels are injected via yticklabels.
% ============================================================
\begin{tikzpicture}
  \begin{groupplot}[
    group style={
      group size=3 by 1,
      horizontal sep=22mm,
    },
    protea plot,
    width=5.5cm, height=6.0cm,
    enlarge y limits=0.06,
    xmin=0,
    xlabel={LightGBM gain (thousands)},
    label style={font=\scriptsize},
    yticklabel style={font=\scriptsize\ttfamily},
    nodes near coords,
    nodes near coords style={font=\tiny, /pgf/number format/fixed, /pgf/number format/precision=0},
    every axis plot/.append style={fill opacity=0.85},
    grid=both,
  ]

    % ── NK panel ────────────────────────────────────────────
    \nextgroupplot[
      title={NK},
      xbar, bar width=3.6pt,
      symbolic y coords={dist,simnw,kpos,nbstd,alnsc,votes,idnw,evcode,gofreq,refdens},
      yticklabels={distance,similarity\_nw,k\_position,nb\_dist\_std,align\_score\_nw,vote\_count,identity\_nw,evidence\_code,go\_term\_freq,ref\_ann\_density},
      ytick=data,
      xmax=1100,
    ]
      \addplot+[fill=proteagreen, draw=proteagreen!60!black] coordinates {
        (971,refdens) (280,gofreq) (163,evcode) (157,idnw) (137,votes)
        (122,alnsc)   (108,nbstd)  (73,kpos)    (70,simnw) (68,dist)
      };

    % ── LK panel ────────────────────────────────────────────
    \nextgroupplot[
      title={LK},
      xbar, bar width=3.6pt,
      symbolic y coords={qual,dist,alnsc,simnw,idnw,nbstd,votes,evcode,gofreq,refdens},
      yticklabels={qualifier,distance,align\_score\_nw,similarity\_nw,identity\_nw,nb\_dist\_std,vote\_count,evidence\_code,go\_term\_freq,ref\_ann\_density},
      ytick=data,
      xmax=350,
    ]
      \addplot+[fill=linkblue, draw=linkblue!60!black] coordinates {
        (309,refdens) (256,gofreq) (169,evcode) (122,votes)  (71,nbstd)
        (64,idnw)     (50,simnw)   (46,dist)    (45,alnsc)   (42,qual)
      };

    % ── PK panel ────────────────────────────────────────────
    \nextgroupplot[
      title={PK},
      xbar, bar width=3.6pt,
      symbolic y coords={taxd,nbstd,votes,kpos,qual,refdens,idnw,simnw,evcode,gofreq},
      yticklabels={tax\_distance,nb\_dist\_std,vote\_count,k\_position,qualifier,ref\_ann\_density,identity\_nw,similarity\_nw,evidence\_code,go\_term\_freq},
      ytick=data,
      xmax=780,
    ]
      \addplot+[fill=BurntOrange, draw=BurntOrange!60!black] coordinates {
        (680,gofreq)  (531,evcode) (307,simnw) (267,idnw)  (262,refdens)
        (185,qual)    (126,kpos)   (117,votes) (88,nbstd)  (75,taxd)
      };
  \end{groupplot}
\end{tikzpicture}

  \caption[Top-10 LightGBM feature importances per category (full-feature, leakage-free)]{%
    Top-10 features by LightGBM gain aggregated over the three per-aspect
    models per category (NK, LK, PK), seed~42, \texttt{bench-v1-K5-filtered}.
    Aggregation features (\texttt{neighbor\_vote\_fraction},
    \texttt{go\_term\_frequency}) dominate across all three categories.
    GO-embedding similarity (\texttt{anc2vec\_neighbor\_cos},
    \texttt{anc2vec\_query\_known\_maxcos}) contributes prominently in LK and PK.
    Raw embedding distance (\texttt{distance}) does not appear in the top-10
    for any category, confirming that the model uses embedding similarity
    primarily as a candidate-generation filter.
    Source: \texttt{figures/fig\_feature\_importance.py} @ \texttt{21d74f7}
    (data: \texttt{figures/data/feature\_importance\_v3.csv}).%
  }
  \label{fig:feature-importance}
\end{figure}

Three observations stand out:
\begin{itemize}
  \item \textbf{Vote-aggregation features dominate}: \texttt{neighbor\_vote\_fraction}
    is the top feature in NK and LK; \texttt{go\_term\_frequency} tops PK. These
    capture the statistical reliability of the annotation vote.
  \item \textbf{GO-embedding features are prominent}: \texttt{anc2vec\_neighbor\_cos}
    and \texttt{anc2vec\_query\_known\_maxcos} appear in all three category
    top-10 lists, reflecting the utility of GO-ontology-derived embeddings as a
    complementary signal to protein-sequence-derived embeddings.
  \item \textbf{Embedding distance ranks low}: the raw cosine distance between
    protein language model embeddings does not appear in any category top-10,
    confirming that the re-ranker learns to use it as a candidate-generation
    filter rather than a final ranking signal.
\end{itemize}

\subsection{Improvement Over Per-Category Configuration (Indicative)}

\Cref{tab:improvement} summarises the absolute $F_{\max}$ change of the
leakage-free full-feature re-ranker relative to the pre-leakage-fix per-category
baseline. These deltas are not clean apples-to-apples comparisons because the
datasets and feature schemas differ; they are reported as an indicative reference.
A controlled leakage-free ablation across all methods on the shared
SELECT-window split is future work.

\begin{table}[htbp]
  \centering
  \caption[Indicative Fmax delta: full-feature (leakage-free) vs.\ per-category (pre-leakage-fix)]{%
    Indicative $F_{\max}$ delta: full-feature leakage-free (\texttt{bench-v1-K5-filtered})
    vs.\ per-category pre-leakage-fix (\texttt{bench-v1-K5}). Datasets and feature
    schemas differ; these figures are not a controlled ablation.}
  \label{tab:improvement}
  \begin{tabular}{lccc}
    \toprule
    \textbf{Category} & \textbf{BPO} & \textbf{MFO} & \textbf{CCO} \\
    \midrule
    NK & +0.174 & +0.039 & +0.086 \\
    LK & +0.134 & +0.031 & +0.087 \\
    PK & $-$0.034 & $-$0.004 & +0.011 \\
    \bottomrule
  \end{tabular}
\end{table}

The large positive deltas in NK-BPO and LK-BPO reflect both the
leakage-fix correction and the expanded feature set (anc2vec). The negative
PK-BPO delta is consistent with the leakage fix tightening the category
criterion, reducing over-optimistic positive rates in PK. A controlled
leakage-free re-run of the per-category configuration is required to isolate the
net feature contribution.

% ── Experiment 7 ──────────────────────────────────────────────
\section{Experiment~7: Comparison with eggNOG-mapper}
\label{subsec:exp7-eggnogmapper}

To position PROTEA relative to an established annotation transfer tool,
the same 20,281-protein test set was submitted to eggNOG-mapper~v2.1.13
\cite{cantalapiedra2021eggnogmapper} using DIAMOND sequence search with
\texttt{--go\_evidence experimental}. Of the 20,281 proteins, 17,334
(85.5\%) received at least one GO term prediction. Full setup and
reproduction instructions are given in \cref{app:external-comparison}.

\begin{table}[htbp]
  \centering
  \caption[Fmax: eggNOG-mapper vs.\ PROTEA methods]{$F_{\max}$: eggNOG-mapper vs.\ PROTEA methods (IA-weighted).}
  \label{tab:results-emapper}
  \resizebox{\textwidth}{!}{%
  \begin{tabular}{l ccc ccc ccc}
    \toprule
    & \multicolumn{3}{c}{\textbf{NK}} & \multicolumn{3}{c}{\textbf{LK}} & \multicolumn{3}{c}{\textbf{PK}} \\
    \cmidrule(lr){2-4} \cmidrule(lr){5-7} \cmidrule(lr){8-10}
    \textbf{Method} & BPO & MFO & CCO & BPO & MFO & CCO & BPO & MFO & CCO \\
    \midrule
    eggNOG-mapper 2.1.13$^\ddagger$ & 0.247 & 0.359 & 0.386 & 0.382 & 0.334 & 0.450 & 0.190 & 0.199 & 0.325 \\
    PROTEA baseline$^\dagger$      & 0.412 & 0.590 & 0.668 & 0.467 & 0.558 & 0.675 & 0.187 & 0.278 & 0.325 \\
    \textbf{PROTEA full-feature} & \textbf{0.599} & \textbf{0.646} & \textbf{0.775} & \textbf{0.620} & \textbf{0.606} & \textbf{0.794} & \textbf{0.165} & \textbf{0.293} & \textbf{0.346} \\
    \bottomrule
  \end{tabular}
  }
\end{table}

\noindent$^\dagger$~Pre-leakage-fix \texttt{bench-v1-K5} dataset (22-feature schema).
$^\ddagger$~eggNOG-mapper evaluated on the same pre-leakage-fix evaluation set
as the baseline; a re-evaluation against the updated split is pending.

The full-feature re-ranker outperforms both eggNOG-mapper and the pre-leakage-fix
baseline in NK and LK across all aspects. In PK-BPO the leakage-fix
recategorisation reduces the positive rate, leading to lower absolute
$F_{\max}$ compared to the pre-leakage-fix baseline and eggNOG-mapper;
this is expected, as the leakage-free criterion is stricter. A side-by-side
comparison on a shared leakage-clean SELECT-window eval set (including a re-run
of eggNOG-mapper) is future work.

The PROTEA baseline surpasses eggNOG-mapper in 8 of 9 cells (pre-leakage-fix
numbers), consistent with the earlier finding that protein language model
embeddings provide a stronger retrieval signal than DIAMOND-based orthology
transfer for NK and LK proteins.

Three caveats apply: (i)~the full-feature and baseline rows use different evaluation
splits and feature schemas and cannot be directly compared;
(ii)~eggNOG-mapper assigns uniform confidence (no threshold optimisation),
which disadvantages it under $F_{\max}$; and (iii)~the eggNOG database
version (5.0.2) predates the evaluation window. A more detailed analysis is
provided in \cref{app:external-comparison}. A clean side-by-side comparison
on the leakage-clean SELECT-window split is future work, covered by the
score-ablation track.

\begin{figure}[htbp]
  \centering
  % ============================================================
%  Figure: PROTEA (baseline + v3) vs eggNOG-mapper 2.1.13
%  Data source: Table 6.6 (chapters/06_evaluation.tex, Experiment 7)
%  NOTE: baseline and eggNOG-mapper rows use pre-leakage-fix bench-v1-K5
%        (22-feature schema, GOA 220->229).
%        v3 uses leakage-free bench-v1-K5-filtered
%        (52-feature schema, GOA 220->230, mean of 3 seeds).
% ============================================================
\begin{tikzpicture}
  \begin{groupplot}[
    group style={
      group size=3 by 1,
      horizontal sep=10mm,
      ylabels at=edge left,
    },
    protea plot,
    width=6.0cm, height=5.0cm,
    enlarge x limits=0.25,
    symbolic x coords={BPO, MFO, CCO},
    xtick=data,
    ymin=0, ymax=0.85,
    ylabel={$F_{\max}$},
    grid=both,
    legend columns=3,
    legend to name=emapper-legend,
    every axis plot/.append style={fill opacity=0.85},
    nodes near coords,
    nodes near coords style={font=\tiny, /pgf/number format/fixed, /pgf/number format/precision=2, /pgf/number format/zerofill},
  ]

    % ── NK panel ────────────────────────────────────────────
    \nextgroupplot[title={NK}, ybar, bar width=4.5pt]
      \addplot+[fill=black!35, draw=black!60]
        coordinates {(BPO,0.247) (MFO,0.359) (CCO,0.386)};
      \addlegendentry{eggNOG-mapper 2.1.13}
      \addplot+[fill=linkblue, draw=linkblue!60!black]
        coordinates {(BPO,0.412) (MFO,0.590) (CCO,0.668)};
      \addlegendentry{PROTEA baseline}
      \addplot+[fill=proteagreen, draw=proteagreen!60!black]
        coordinates {(BPO,0.599) (MFO,0.646) (CCO,0.775)};
      \addlegendentry{PROTEA re-ranker v3}

    % ── LK panel ────────────────────────────────────────────
    \nextgroupplot[title={LK}, ybar, bar width=4.5pt]
      \addplot+[fill=black!35, draw=black!60]
        coordinates {(BPO,0.382) (MFO,0.334) (CCO,0.450)};
      \addplot+[fill=linkblue, draw=linkblue!60!black]
        coordinates {(BPO,0.467) (MFO,0.558) (CCO,0.675)};
      \addplot+[fill=proteagreen, draw=proteagreen!60!black]
        coordinates {(BPO,0.620) (MFO,0.606) (CCO,0.794)};

    % ── PK panel ────────────────────────────────────────────
    \nextgroupplot[title={PK}, ymax=0.42, ybar, bar width=4.5pt]
      \addplot+[fill=black!35, draw=black!60]
        coordinates {(BPO,0.190) (MFO,0.199) (CCO,0.325)};
      \addplot+[fill=linkblue, draw=linkblue!60!black]
        coordinates {(BPO,0.187) (MFO,0.278) (CCO,0.325)};
      \addplot+[fill=proteagreen, draw=proteagreen!60!black]
        coordinates {(BPO,0.165) (MFO,0.293) (CCO,0.346)};
  \end{groupplot}

  \node[anchor=north] at ($(group c2r1.south) + (0,-7mm)$)
    {\pgfplotslegendfromname{emapper-legend}};
\end{tikzpicture}

  \caption[PROTEA vs eggNOG-mapper across NK/LK/PK and aspects]{%
    Head-to-head $F_{\max}$ comparison between eggNOG-mapper~v2.1.13,
    PROTEA's pre-leakage-fix embedding-only baseline ($^\dagger$), and
    PROTEA's leakage-free full-feature re-ranker on the GOA~220$\to$230
    evaluation set. The full-feature and baseline rows use different datasets;
    direct cell-by-cell comparisons are indicative only. NK and LK show large
    margins over eggNOG-mapper across all methods. PK-BPO shows lower full-feature
    performance due to the leakage-fix recategorisation tightening the positive set.
    Source: \texttt{figures/fig\_protea\_vs\_emapper.py} @ \texttt{21d74f7}
    (data: \texttt{figures/data/protea\_vs\_emapper.csv}).%
  }
  \label{fig:protea-vs-emapper}
\end{figure}
